\section{Antecedentes y plan de trabajo}
\label{anx:antecedentes}
% Origen: 01-introduccion.tex (antecedentes) y 02-paso1-formulacion.tex (plan de trabajo), movidos al anexo.

\subsection{Antecedentes}

Varios estudios recientes simularon estaciones de metro con AnyLogic y midieron parámetros que este
trabajo usa como referencia.
\begin{itemize}
  \item \textcite{seriani2025} observaron los molinetes de una estación del Metro de Valparaíso y midieron
        un tiempo de paso medio de 2~s. Registraron además un sesgo fuerte por cercanía: el molinete
        más próximo al recorrido es el más usado. Validaron su modelo comparando el uso simulado de cada
        molinete con el observado.
  \item \textcite{wang2024} midieron con video 813 pasos por molinete en Pekín, con medias de 2,91 a 3,49~s
        según el medio de pago.
  \item \textcite{tian2026} modelaron una estación de trasbordo que recibe pasajeros de un ferrocarril, la
        situación análoga a la del Roca y la Línea~C. Encontraron que el 88,2~\% usa el grupo de molinetes
        más cercano a la salida del tren.
  \item \textcite{li2024} calibraron con video el modelo de fuerza social de AnyLogic y obtuvieron una
        velocidad deseada de 1,37~m/s, dentro del rango de 1,1 a 1,5~m/s de la literatura.
  \item \textcite{zhen2024} usaron datos abiertos de molinetes de Nueva York, como aquí los de SBASE, pero
        con arribos de Poisson y sin oleadas.
\end{itemize}
Lo que agrega este trabajo es representar las oleadas a partir del horario real de los trenes y
distribuir entre ellos la demanda observada por ventana.

\subsection{Plan de trabajo}
\label{anx:plan}

El estudio se organizó en fases, cada una con un entregable verificable (tabla~\ref{tab:plan}).

\begin{table}[H]
\centering
\caption{Plan de trabajo}
\label{tab:plan}
\small
\begin{tabularx}{\textwidth}{@{}l X l@{}}
\toprule
Fase & Contenido & Fecha \\
\midrule
Definición & Aprobación del tema, perfilado del dataset de SBASE y definición del caso & 20 a 23/09/2026 \\
Pedidos de datos & Correos a SBASE, Emova y Trenes Argentinos; pedidos formales por Ley~104 y Ley~\num{27275} & 23/09 y 05/10/2026 \\
Modelo & Capa lógica, capa peatonal, demanda por horario del Roca y automatización de las corridas & 20/09 a 24/09/2026 \\
Pilotos & Diagnóstico del atasco, corrección y calibración del escenario intermedio & 24/09 a 05/10/2026 \\
Producción inicial & 30 pares \EO{}-\EI{} con un día promedio y reproducción en una segunda máquina & 05/10 y 09/10/2026 \\
Modelo con días reales & Días hábiles y trenes reales, accesos por las escaleras del bloque central, \ET{} y \EIT{} & 02/11 a 09/11/2026 \\
Producción final & 6 días hábiles reales $\times$ 4 escenarios, en paralelo (24 corridas) & 10 y 11/11/2026 \\
Informe y video & Redacción en \LaTeX{}, presentación, grabación y defensa & 10 a 11/11/2026 \\
\bottomrule
\end{tabularx}
\end{table}
